\section{Discussion and Conclusions}
\label{sec:discussion}

We can now state the main differences between \ourapproach and the
two approaches that inspired it, against the background of what the rest of the paper has
shown. The main lesson from the literature on actual causality is that intuitive causal
attribution often requires holding some endogenous variables, members of a ``witness
set'', fixed at their factual values. This ensures context sensitivity by controlling
complicating causal paths. We treat witnesses differently: instead of looking for
a single witness set that reveals a causal relationship, we use a distribution over
witness sets and integrate the causal impact score across them. As in many probabilistic
generalisations, this focuses on whole distributions, and it lets
\ourapproach account for the number of witness sets and the strength of
within-witness-set effects. Moreover, the construction admits off-the-shelf probabilistic
inference in place of worst-case subset enumeration. The relationship between the two
views is now formal where it was once gestural: under matching choices of $\Gamma$ and
$\Delta$, AC verdicts coincide with the support of $\sum_{\mathbf{T}}\Phi$
(Theorems~\ref{th:ac-exp}--\ref{th:exp-ac} and Proposition~\ref{prop:filter}), and the
empirical companion in \cref{sec:ac_benchmark} confirms this at scales where AC
enumeration is intractable.

The main concept borrowed from Pearl's probability of necessity and sufficiency is the
joint distribution over outcomes from two interventional regimes, pushed forward through
a user-chosen causal impact function. The key difference is that we use witness sets to
allow for context-sensitivity and to break symmetries present in otherwise pathological
cases like overdetermination and undercutting. Section~\ref{sec:shap_examples} works
through the binary instantiation in which the necessity and sufficiency components
coincide with PN and PS, and \cref{sub:synthetic} reads the two components apart on
a continuous synthetic model where each is needed to bring out a different causal
archetype.

The further generalisations are mostly bookkeeping in the same spirit. Instead of fixing
a single alternative value for a potential cause, we integrate over a factual-excised
alternative-value distribution; this generalises immediately to continuous variables.
Instead of asking whether a particular set is an actual cause, we sample suspect sets and
gauge a feature's responsibility in terms of its participation in sets with causal
impact. Instead of asking whether the outcome differs, a customisable causal impact
function $ci$ measures the differences across counterfactual worlds, including between
continuous variables. The variable-selection distribution $\Gamma$, the alternative-value
distribution $\Delta$, and the impact function $ci$ become explicit, user-facing
components of the framework, no longer implicit defaults.

Against SHAP-family attributions, the framework's central commitment is to
\emph{intervention with witnesses} over \emph{observational marginalisation}. The
controlled comparisons of Section~\ref{sec:shap_examples} show that this commitment does
real work on canonical overdetermination and continuous-mediation examples, where SHAP
and Causal SHAP fail in characteristic ways that \ourapproach avoids. The synthetic
benchmark of \cref{sub:synthetic} confirms this on a model with known ground
truth, and the case study of \cref{sec:avm} reports the same qualitative pattern
on a production-grade trained model with millions of data points. A separate comparison
to gradient-based attribution, where the difference is one of differentiability and locality, not the
causal/non-causal axis, is reported in \cref{sec:DCE}. The
thread connecting these results is structural: the witness mechanism that satisfies
D-A-rank for Alice in the binary OBCB walkthrough of
Section~\ref{sec:shap_examples} is the same machinery driving the upstream-versus-downstream
attribution split on the AVM in \cref{sec:avm}.

A broader test is how \ourapproach fares against the desiderata that the XAI
literature typically asks counterfactual-explanation methods to meet, of which
\cite{verma2022recourses} list five. \emph{Sparsity}, the requirement that
explanations change as few features as possible, follows directly from the
cardinality cap on suspect sets in $\Gamma$, which biases the search toward small
$\mathbf{C}$. \emph{Multiple counterfactuals} (the ability to return several
alternative explanations instead of committing to one) falls out of the Monte Carlo
construction, since a single sweep over $\Gamma$ and $\Delta$ produces a population of
$(\mathbf{C}, \mathbf{c}', \mathbf{T})$ triples, not a single optimum.
\emph{Causal-constraint satisfaction}, the demand that explanations respect
downstream propagation rather than treat features as independently editable, is
automatic once suspect interventions are routed through the causal model: changing
$\mathbf{C}$ propagates to descendants by construction. \emph{Model-agnosticism}
(applicability to arbitrary predictors) \ourapproach satisfies only partially: it
can wrap a black-box predictor, but at the cost of the causal sensitivity that
motivates the framework in the first place.

\emph{Actionability}, the requirement that explanations not recommend changes to
features users cannot in practice change, is where we part company with the
literature. The standard recommendation is to fold ease-of-change into the
cause-identification procedure itself, weighting features by mutability so that
immutable ones are never returned as causes (\cite{verma2022recourses}; cf.\
\cite{Karimi2021}, who shift the focus from explanations altogether to interventions).
We think this conflates two distinct questions. A procedure that suppresses gender as
a possible cause whenever a loan decision is at stake will help applicants satisfy the
existing system without surfacing the kind of structural
problem that responsibility attribution can in principle reveal. Cost-of-action
layers can sit downstream of \ourapproach for users who need actionable
recommendations, but they should compose with responsibility attribution instead of
replacing it.

Several limitations are worth flagging. \Ourapproach presumes access to a trained
probabilistic causal model from which counterfactual outcomes can be sampled, a
stronger requirement than methods that work from observational data and a causal graph
alone. The Monte Carlo variance of the estimator grows with model size and with the
cardinalities of suspect and witness sets; \cref{sec:ac_benchmark} reports
empirical scaling and an ablation that isolates the contribution of witness pinning, but
sample-efficient estimators for high-cardinality settings remain an open direction.
Finally, the AVM case study is deliberately limited in disclosure: the apples-to-apples
SHAP-vs-PCI comparison rests on the synthetic benchmark of \cref{sub:synthetic},
with the production model serving as a feasibility witness, not as quantitative
evidence.



In sum, \ourapproach provides a scalable, context-sensitive causal explanation
method for causal probabilistic machine-learning models. A natural probabilistic
generalisation of ideas already in the literature leads to a tractable estimator that can
be paired with standard inference algorithms. We benchmarked the approach against exact
actual-causality computations, characterised its qualitative behaviour on synthetic
models with known ground truth and on canonical SHAP failure cases, and showed its
computational feasibility on a production-grade model.

Prior efforts to mechanise actual-causality reasoning in software --- notably the Actual
Causality Canvas of \citet{ibrahimCanvas2020} and the SAT- and optimization-based
computations of \citet{ibrahimChecking2020}, which operationalise the Halpern--Pearl
definition for forensic investigation, fault diagnosis, and explainable AI --- are built
for \emph{discrete} (typically binary) structural models, where the actual cause is
recovered by combinatorial search over candidate sets. The complementary regime of large,
continuous-valued models has received less attention, and it is there that the exact
notions become harder to apply directly. \Ourapproach develops approximate causal
attribution workable for users of models of non-trivial size, with the design choices left
to the user.
